\documentclass[11pt,letterpaper]{article}
\usepackage[margin=0.9in]{geometry}
\usepackage[T1]{fontenc}
\usepackage{lmodern}
\usepackage{microtype}
\usepackage{array,longtable,booktabs}
\usepackage{xcolor}
\usepackage[hidelinks]{hyperref}
\usepackage{enumitem}
\setlist[itemize]{leftmargin=*,nosep}
\setlength{\parindent}{0pt}
\setlength{\parskip}{0.5em}
\newcommand{\code}[1]{\texttt{\detokenize{#1}}}
\title{HW2: Code-Action and MCP Agents\\\large Design Report}
\author{Amit Bal \\ \texttt{asb2332}}
\date{}

\begin{document}
\maketitle

\section{MCP Tool Contracts and Call Path}

I expose analysis, source access, editing, and rescanning through a single stdio MCP server. The model selects from the tools discovered from that server; it does not call the server's Python functions directly. Each contract specifies the arguments, the returned information, and the conditions under which a request is refused. Paths are relative to the workspace; line numbers are one-based.

\subsection{Model-facing tool contracts}

In the table, optional arguments are marked with a question mark. The result column summarizes the data returned to the host; the MCP SDK wraps scalar and list results in a \code{result} field in structured output where an output schema is available.

{\small
\renewcommand{\arraystretch}{1.15}
\begin{longtable}{@{}p{0.205\textwidth}p{0.285\textwidth}p{0.46\textwidth}@{}}
\toprule
\textbf{Tool} & \textbf{Arguments} & \textbf{Result and purpose} \\
\midrule
\endfirsthead
\toprule
\textbf{Tool} & \textbf{Arguments} & \textbf{Result and purpose} \\
\midrule
\endhead
\bottomrule
\endfoot
\code{locate} & \code{path}: string, \code{line}: integer & Enclosing callable signature, or \code{module scope}; locates the code around an alert. \\
\code{source} & \code{signature}: string & Source text for the requested callable. \\
\code{callers} & \code{name}: string, \code{module}?: string or null & List of callers and their locations. \\
\code{callees} & \code{name}: string, \code{module}?: string or null & List of callees, including external targets. \\
\code{dataflow} & \code{path}: string, \code{line}: integer & The first page of data-dependence edges for the enclosing callable, with \code{edges}, \code{total}, and \code{complete}. \\
\code{entry_points} & None & Detected route handlers and CLDK's entry-point coverage report. \\
\code{read_file} & \code{path}: string; \code{start}?, \code{end}?: integer or null & File text, optionally restricted to an inclusive line range. \\
\code{search} & \code{pattern}: string; \code{path}?: string or null & Up to ten matching Python-source lines, each with a file, line number, and text. \\
\code{edit} & \code{path}, \code{old}, \code{new}: strings & Replaces one exact occurrence in the selected file and returns a patch status, diff, and acceptance checks. \\
\code{rescan} & \code{path}: string & Bandit alerts for the selected file after an edit. \\
\end{longtable}
}

The server refuses a file request outside its allowed workspace or current alert scope. On OWASP, file reads are restricted to the selected case; edits and rescans are restricted to the selected alert's file. For example, \code{edit} refuses an empty \code{old} value or one that does not occur exactly once. \code{read_file} refuses invalid line ranges, and \code{search} refuses an invalid regular expression. Analysis requests can also be refused when a file or callable is unavailable or outside the selected scope. The host selects the alert using \code{set_alert}; it does not offer that control, \code{configure_alerts}, or \code{validate_evidence} to the model.

For OWASP, I pass the selected file to CLDK as the analysis target and filter location-bearing results to the selected case where applicable. Consequently, a selected-file query is not equivalent to a whole-project query: a helper may appear as an external callee, and the entry-point list excludes handlers from other cases.

\subsection{A \code{read_file} request through MCP}

\begin{enumerate}[leftmargin=*,itemsep=0.25em]
\item \textbf{Discovery.} The host starts the configured stdio server, initializes an MCP session, and calls \code{tools/list}. It removes the three host-only controls, records the offered names in the discovery trace event, and binds the remaining definitions to the model.
\item \textbf{Model and dispatcher.} The model requests \code{read_file} with a path and optional range. The graph passes the request to the dispatcher, which records it and checks the name and arguments against the discovered input schema. It refuses an unknown, host-only, or malformed request rather than forwarding it.
\item \textbf{Client and server.} For an allowed request, the MCP client sends \code{tools/call} as JSON-RPC over the server process's standard input and output. The server checks the selected-case restriction and calls its workspace read helper. The helper resolves the path within the workspace and reads the requested text.
\item \textbf{Return to the model.} The server's result or tool error travels back through the client. The dispatcher records the adjacent result event, including the refusal layer when applicable, and the graph adds the observation to the conversation for the next model call.
\end{enumerate}

\subsection{What the contract enforces}

The server's tool descriptions explain intended behavior to the model; their wording is not an access-control mechanism. The dispatcher checks the discovered argument schema, and the MCP SDK validates tool arguments and any declared output schema. Those shape checks do not establish that a path is authorized or that returned facts are correct. The server's path and selected-alert checks enforce the request-specific restrictions, while the sandbox supplies the separate operating-system boundary.

All registered tools, including the host-only controls, expose generated output schemas. Typed result definitions specify the fields of structured responses, including data-flow edges, entry points, and patch checks. The MCP SDK validates returned results against these schemas; malformed results are returned as tool errors. This validation checks the shape of a result, not the correctness of its contents.

\section{Access Rules and Sandbox Policy}

I apply the same intended read, write, and network rules to both agents. They may read only target code, never \code{.git} or the answer key; on OWASP, they may read only the case for the current alert. During remediation, they may edit only the file associated with that alert, inside the disposable workspace, and must not execute target code. Network access is limited to the model provider.

The fixed restrictions apply throughout a run: the answer key and \code{.git} are withheld, target writes are limited to workspace copies, and network access is limited to the provider. The current OWASP case and the file that may be edited are per-alert restrictions: their allowed values change as the runner moves between alerts.

\subsection{Where the restrictions are enforced}

\textbf{Prompt.} Both prompts state the intended restrictions, including the rule against executing target code. The MCP prompt additionally directs the model to treat source code as untrusted. These instructions can be bypassed if the model follows malicious text in a source file, so they are not the enforcement boundary.

\textbf{MCP server.} The host calls \code{set_alert} for each alert and excludes that control from the model's offered tools. The server uses the selected file to restrict OWASP reads and edits; its path resolver also rejects absolute paths, parent traversal, and paths resolving outside the workspace. A forbidden request that reaches the server is returned as an MCP tool error; the dispatcher records whether the refusal came from the server or from a kernel permission error. These checks constrain requests through the server, not arbitrary code running in the server's own process: a sandbox process that directly opens a file bypasses the MCP checks.

\textbf{CodeAct.} The harness creates a workspace and a persistent Python namespace for an alert. The model's generated script can read or modify files directly; its internal operations do not pass through an MCP server for per-operation checks. CodeAct also retains native dispatcher tools with their own checks, but a script can bypass those tools. The prompt states the per-alert rules and the sandbox enforces its fixed policy; the fixed policy does not know which alert is currently open. Staging and withholding reduce the files in the workspace, but they are not a substitute for an OS-enforced per-alert rule.

\textbf{Operating-system kernel.} Both agents, and the MCP server when used, run under the OpenShell policy. Landlock restricts filesystem access independently of the model's instructions and server tool descriptions. The kernel's fixed grants do not change on \code{set_alert}, so a direct read of another OWASP case can pass the kernel even though the MCP server must refuse it. Conversely, if a request passes the server but the kernel denies the underlying file access, the server reports a kernel refusal. A broad filesystem or network grant would weaken this layer; the kernel can only enforce the policy it is given.

\subsection{Filesystem grants and withheld paths}

The policy enables \code{include_workdir} and grants read-only access to the following paths:
\begin{itemize}
\item \code{/bin}, \code{/usr}, \code{/lib}, and \code{/lib64} provide programs, the Python interpreter, installed packages, and runtime libraries. Paths absent from an image are skipped.
\item \code{/etc} and \code{/proc} provide system configuration and process information; \code{/dev/urandom} supplies randomness. \code{/app} contains the agent, server, dependencies, and configuration.
\item The permitted OWASP entries are \code{/target/testcode}, \code{/target/helpers}, and \code{/target/app.py}. The other target grants are \code{/target/radicale}, the SGLang \nolinkurl{python/sglang/multimodal_gen/runtime} directory, and \nolinkurl{scripts/playground/replay_request_dump.py} under \code{/target}. Each image places one target at \code{/target}; absent paths are skipped.
\end{itemize}
There is no grant on the target root itself. In particular, the root-level answer key and \code{/target/.git} are not granted. The permitted target entries are staged into \code{/sandbox/work/target}, excluding the answer key and \code{.git}. CLDK analyzes a writable workspace copy rather than attempting to create its cache under the read-only target source.

The read-write grants are \code{/sandbox}, for the staged input and outputs; \code{/tmp}, for temporary workspaces, analysis files, and scan output; and \code{/dev/null}, for discarded output. These are broader than a single per-alert directory: the server and prompt still carry the per-alert edit rule. With \code{include_workdir}, the sandbox's working directory is also included in the filesystem policy. The effective policy inspected for the OWASP sandbox additionally listed \code{/var/log} as read-only. Landlock compatibility is \code{hard_requirement}, so startup should fail rather than silently continue without the required enforcement.

\subsection{Network and provider attachment}

The \code{model_provider} policy permits REST access to \code{api.openai.com:443} for \nolinkurl{/usr/local/bin/python3.12}, with \code{enforcement: enforce} and \code{access: read-write}; the latter permits the model API's requests and responses. The OpenAI provider attachment adds an effective \code{_provider_openai} policy for the same endpoint and binary. The provider attachment supplies the credential outside the sandbox; the sandbox receives a placeholder rather than the actual key. Authorizing the interpreter also authorizes its children, so the binary grant is a process-level permission, not a promise that only the intended model call can use it.

\subsection{Observed boundary checks}

In a fresh OWASP sandbox, a direct open of \code{/target/expectedresults-0.1.csv} was denied, as was a write under \code{/target}; a direct read of another OWASP case succeeded. Separately, raw MCP requests to read the answer key or a different case, or to edit a different file, returned tool errors. These are observations for the tested paths, not proof that every possible request is blocked. The planted-injection runs are a separate measurement of model behavior.

\end{document}
